% =====================================================================
\chapter{RELATED WORKS}
\label{ch:related}
% =====================================================================

This chapter first introduces the theoretical foundations on which the
thesis builds (Section~\ref{sec:overview-foundations}): crowdsourced
last-mile delivery, pickup-and-delivery routing, mechanism design, and
column reduction. Section~\ref{sec:litreview} reviews five streams of
literature, classifies them and positions the thesis against the closest
studies. Section~\ref{sec:candidates} describes the candidate methods for
each component of the system, which Chapter~\ref{ch:method} compares.

\section{Overview}
\label{sec:overview-foundations}

\subsection{Crowdsourced last-mile delivery}

\citet{archetti2016vrpod} introduced the \emph{vehicle routing problem with
occasional drivers} (VRPOD): a company that operates its own fleet may also
use ordinary people who are driving to their own destination and who are
willing to make one delivery on the way in exchange for a compensation set by
the company. Later work extended the model to time windows, multiple
deliveries per driver and pickup-and-delivery settings, and studied
platforms in which the crowd is the main source of capacity
\citep{alnaggar2021crowdsourced,savelsbergh2022challenges}. Two features of
crowd capacity recur in this literature and drive the present thesis.
First, crowd drivers are not employees: whether and at what price they
serve is their own decision, so their costs are private. Second, crowd
drivers differ in how they move. \citet{luy2024workforce} distinguish fixed
drivers, gigworkers and occasional drivers in a strategic workforce-planning
model; this thesis adopts that taxonomy (Table~\ref{tab:classes}) at the
operational level, where a single dispatch session is cleared.

\subsection{Pickup-and-delivery routing with time windows}

In a pickup-and-delivery problem with time windows (PDPTW), each request
consists of a pickup and a delivery that must be served by the same vehicle,
the pickup before the delivery, within given time windows and capacity
limits \citep{dumas1991pickup}. A vehicle that arrives before a window
opens waits. Even checking whether a single vehicle can serve a given set of
requests requires searching over visiting sequences; for a bundle of $k$
orders there are $(2k)!/2^k$ precedence-feasible sequences.

Large routing problems are usually formulated as \emph{set-partitioning}
problems: each column is a feasible route, and the model selects columns so
that each request is covered exactly once. The columns are generated by
\emph{column generation}, whose pricing subproblem is an elementary shortest
path problem with resource constraints (ESPPRC), solved by \emph{label-setting}
algorithms \citep{lubbecke2005selected,feillet2004exact}. A label represents
a partial path together with its consumed resources (time, load, cost), and
a label can be discarded if another label at the same node is at least as
good in every resource, so that every extension of the first is matched by a
no-worse extension of the second. For pickup and delivery the classical
dominance rule is that of \citet{dumas1991pickup}, presented in modern form
by \citet{gschwind2018bidirectional}. Section~\ref{sec:candidates} explains
why column generation itself cannot be used in this thesis, although its
labeling machinery can.

\subsection{Mechanism design and the VCG mechanism}

In a procurement mechanism the platform (the buyer) asks each agent $i$ to
report her private \emph{type} $\theta_i$, here her value of time, and then
applies an \emph{allocation rule} that decides who serves what and a
\emph{payment rule} that decides how much each agent is paid. Agent $i$'s
utility is her payment minus her true cost. The mechanism is
\emph{dominant-strategy incentive compatible} (DSIC, or truthful) if
reporting $b_i=\theta_i$ maximises every agent's utility whatever the others
report, and \emph{individually rational} (IR) if a truthful agent never
obtains negative utility \citep{nisan2007agt}.

The Vickrey--Clarke--Groves (VCG) mechanism chooses the allocation that
minimises total reported cost and pays each winner $i$
\[
p_i = c_i(b_i) + V_{-i} - V,
\]
where $V$ is the optimal total cost and $V_{-i}$ the optimal total cost when
$i$ is absent (the Clarke pivot). The difference $V_{-i}-V$ is the
externality that $i$'s presence creates, and $p_i$ does not depend on $i$'s
own report as long as she wins the same bundle, which is why misreporting
cannot help. VCG requires the allocation to be \emph{exactly} optimal: with a
heuristic allocation the truthfulness argument breaks down
\citep{nisan2007computationally}. \citet{nisan2007computationally} also
showed that VCG remains truthful if the optimisation is exact over a
\emph{restricted range} of allocations, provided that range is fixed
independently of the reports (``maximal-in-range'' mechanisms). This is the
property that the route pool of this thesis must satisfy. Computing VCG
payments requires one optimisation without each winner in addition to the
base problem, which \citet{nisan2000computationally} identified as a central
obstacle to deploying VCG in practice.

\subsection{Column dominance and redundancy}

Deleting columns that cannot be needed is a classical preprocessing step. In
set covering, a column is \emph{dominated} if a cheaper column covers a
superset of its rows \citep{beasley1987algorithm}. In column generation,
\citet{vanderbeck1994thesis} calls a column dominated if another column has
no larger reduced cost for all values of the dual variables within their
domains, and \citet{sol1994column} strengthens dominance to
\emph{redundancy}, in which a column is replaced by a combination of its
subcolumns of no smaller total cost (both as summarised by
\citealp{lubbecke2005selected}). Restrictions on the dual space, such as the
dual-optimal inequalities of \citet{valerio2005extra} and
\citet{benamor2006dual}, enlarge the set of columns that can be shown to be
dominated. The local frontier of this thesis turns out to be a form of
redundancy in which the outside option plays the role of such a dual
restriction (Remark~\ref{rem:lp}).

\section{Literature Review}
\label{sec:litreview}

\subsection{Crowdshipping operations}

A large part of the crowdshipping literature optimises routing and
assignment when the crowd's behaviour is given. \citet{archetti2016vrpod}
let occasional drivers serve at most one customer each, for a compensation
set by the company. \citet{torres2022crowdshipping} model crowdshipping as an
open vehicle routing problem with stochastic destinations, and
\citet{torres2022stochastic} study vehicle routing with a stochastic supply
of crowd vehicles and time windows. \citet{mousavi2022stochastic} combine
crowd-shippers with mobile depots in a two-stage stochastic programme in
which crowd-shipper availability is uncertain. \citet{ausseil2022supplier}
let a peer-to-peer platform offer each independent supplier a menu of
requests, trading off acceptance probability against duplicate selections.
\citet{mancini2022bundle} and \citet{mancini2024bundle} generate bundles for
occasional drivers along travel corridors and in a spatial--temporal space,
and \citet{zhou2026service} develop an exact branch-price-and-cut algorithm
for a service-centric vehicle routing problem with crowdshipping.
\citet{luy2024workforce} plan the size of a hybrid workforce of fixed
drivers, gigworkers and occasional drivers over a long horizon. Across this
stream, driver costs or compensations are known inputs or behavioural
response models; none elicits private costs through a mechanism.

\subsection{Procurement mechanisms and auctions for crowdshipping}

A second stream elicits costs. \citet{triki2021combinatorial} lets
occasional drivers bid on bundles alongside a company fleet but does not
establish incentive compatibility. \citet{chen2023truthful} obtain a
truthful reverse auction for occasional couriers in which the couriers
select their own bundles, so bundle selection lies outside the incentive
analysis. \citet{zou2022mechanisms} give an exact VCG mechanism with
Clarke-pivot payments on a platform-generated job pool, with a
non-strategic backup vehicle and depot-based jobs. \citet{oyama2025market}
study market-based matching with task bundling, using a truthful auction
inside a subproblem and a fluid approximation at the master level.
\citet{xu2026incentive} design incentive-compatible auctions for a hybrid
fleet in which only the crowd bids. Closest to this thesis,
\citet{li2026auction} generate a route pool by solving pickup-and-delivery
problems with recursive branching that exploits the spatial monotonicity of
detour time, run exact VCG on it, and propose a greedy second-best mechanism
with a regret bound; their crowd carriers form a single detour-constrained
class.

\subsection{Safe reduction of column pools}

Beyond the classical dominance rules of Section~\ref{sec:overview-foundations},
\citet{muller1998solving} deletes a set-partitioning column when other
columns partition it at lower total cost. Both this rule and that of
\citet{beasley1987algorithm} assume known costs and compare a column against
the whole pool. \citet{pereira2013robust} treat set-covering costs as
intervals, but seek a single min--max-regret solution rather than a
deletion rule that is exact for every cost realisation. On the mechanism
side, \citet{aggarwal2006knapsack} and \citet{kempe2010frugal} prune agents
before running a VCG-type mechanism in order to bound its frugality ratio;
their pruning may depend on bids and is not required to preserve the
optimal value exactly.

\subsection{Label dominance and rollback pruning}

Label-setting algorithms for the ESPPRC rely on dominance between labels at
the same node \citep{feillet2004exact}. \citet{lozano2016exact} add
\emph{rollback pruning}, which compares a partial path with the same path
skipping a node: if skipping the node is cheaper and no less feasible, the
path through it is discarded. The argument relies on two properties of the
ESPPRC: costs are additive over arcs, and arriving earlier at a node never
costs anything. In the procurement setting of this thesis neither holds
exactly, because working time is priced and waiting is allowed; this is why
the rule of Section~\ref{sec:sd-label} needs a correction.

\subsection{Computing VCG payments}

\citet{hershberger2001vickrey} show that for shortest-path auctions all
VCG payments can be obtained far faster than by independent re-solves.
% VERIFY: exact statement of Hershberger & Suri (2001) before submission.
\citet{bleischwitz2005accelerating} share bounds across the simultaneous
searches of a combinatorial auction; \citet{sandholm2005cabob} decompose a
single winner-determination solve along the components of a bid-conflict
graph; and \citet{bunz2015faster} prove a separability result for
core-constraint generation based on a bipartite conflict graph rebuilt at
every iteration. \citet{wilkens2012single} and \citet{li2026auction} avoid
the repeated solves altogether by relaxing exactness, at the cost of
truthfulness only in expectation or of an approximate allocation.

\subsection{Summary, classification and positioning}

Table~\ref{tab:litclass} classifies the reviewed approaches by stream and by
what they contribute to the problem of this thesis, and
Table~\ref{tab:positioning} compares the closest studies attribute by
attribute.

\begin{table}[t]
\centering
\caption{Classification of the reviewed literature}
\label{tab:litclass}
\small
\begin{tabular}{P{2.6cm}P{4.2cm}P{3.4cm}P{3.6cm}}
\toprule
Stream & Representative studies & What it provides & What is missing for this thesis\\
\midrule
Crowdshipping operations & \citet{archetti2016vrpod}; \citet{torres2022crowdshipping,torres2022stochastic}; \citet{mousavi2022stochastic}; \citet{ausseil2022supplier}; \citet{mancini2022bundle,mancini2024bundle}; \citet{zhou2026service}; \citet{luy2024workforce} & Routing models for occasional drivers and hybrid fleets; bundle generation & Costs are known inputs; no elicitation of private costs\\
\addlinespace
Procurement mechanisms & \citet{triki2021combinatorial}; \citet{chen2023truthful}; \citet{zou2022mechanisms}; \citet{oyama2025market}; \citet{xu2026incentive}; \citet{li2026auction} & Auctions and truthful payments for crowd carriers & One strategic class only; no characterisation of pool reduction\\
\addlinespace
Column pool reduction & \citet{beasley1987algorithm}; \citet{muller1998solving}; \citet{vanderbeck1994thesis}; \citet{sol1994column}; \citet{pereira2013robust} & Dominance and redundancy of columns & Known costs or LP duals; no report domain; no minimality result\\
\addlinespace
Label dominance & \citet{dumas1991pickup}; \citet{feillet2004exact}; \citet{gschwind2018bidirectional}; \citet{lozano2016exact} & Safe pruning of partial paths & Assumes unpriced time; unsafe with priced waiting\\
\addlinespace
VCG payment computation & \citet{nisan2000computationally}; \citet{hershberger2001vickrey}; \citet{sandholm2005cabob}; \citet{bunz2015faster}; \citet{wilkens2012single} & Faster or approximate counterfactuals & No decomposition over a static conflict graph with a priced outside option\\
\bottomrule
\end{tabular}
\end{table}

\begin{table}[t]
\centering
\caption{Positioning against the closest related studies}
\label{tab:positioning}
\footnotesize
\setlength{\tabcolsep}{3pt}
\begin{tabular}{P{2.4cm}P{1.6cm}P{1.7cm}P{1.5cm}P{2.2cm}P{2.0cm}P{2.0cm}}
\toprule
Study & Open-route class bids & Detour class bids & Outside option & Bundle source & Incentive property & Pool reduction characterised \\
\midrule
\citet{archetti2016vrpod} & Company fleet & No bidding & -- & -- & -- & No \\
\citet{triki2021combinatorial} & Company fleet & Yes & -- & Bids & Not shown & No \\
\citet{zou2022mechanisms} & No & Yes (depot-based) & Backup vehicle & Platform & Exact VCG & No \\
\citet{mancini2022bundle} & No & -- & -- & Corridors & -- & No \\
\citet{chen2023truthful} & No & Yes & -- & Couriers & DSIC, fixed bundle & No \\
\citet{luy2024workforce} & Known cost & Known cost & Fixed drivers & -- & -- & No \\
\citet{oyama2025market} & One population & -- & Opt-out & Task chains & Truthful sub-auction & No \\
\citet{xu2026incentive} & Non-strategic & Yes & Yes & Pre\-determined & Yes & No \\
\citet{li2026auction} & No & Yes & Fixed price & Platform & Exact VCG, greedy & Detour-monotone branching \\
This thesis & Yes & Yes & Fixed price & Platform, bid-independent & Exact VCG & Frontier and tightness \\
\bottomrule
\end{tabular}
\par\smallskip\footnotesize Note: each row should be re-checked against the primary source before final submission.
\end{table}

\paragraph{Similarities and differences with the key references.}
The thesis shares its mechanism with \citet{zou2022mechanisms} and
\citet{li2026auction}: exact VCG with Clarke-pivot payments on a
platform-generated pool. It differs from both in having two strategic crowd
classes with different route geometries, and from \citet{li2026auction} in
that their pool reduction exploits a detour budget, which gigworkers do not
have. It shares its driver taxonomy with \citet{luy2024workforce}, but
treats costs as private and works at the level of one dispatch session. Its
frontier is a form of the redundancy of \citet{sol1994column} in which the
substitute consists of exactly one subroute of the same driver together
with outside-option columns. The restriction matters: Sol shows that no
column is redundant when all subproblems differ from one another, and here
every driver is a distinct subproblem; it is the publicly priced outside
option, whose columns belong to no driver, that restores redundancy.
% VERIFY: Sol's statement is read from Lubbecke & Desrosiers (2005); check Sol (1994).
Two further features have no counterpart in these sources: costs depend
on a report ranging over a domain fixed by the mechanism designer, and the
substitute may vary with the report; and safety is required for the integer
optimum and for every counterfactual problem, not for pricing a linear
relaxation. Whether the frontier is also minimal over all local
bid-independent rules is planned for the final thesis. The label rule
shares the idea of rollback pruning with \citet{lozano2016exact}, but must
correct for waiting. The payment decomposition uses a unipartite conflict
graph built once from the pool, unlike the per-solve decomposition of
\citet{sandholm2005cabob} and the per-iteration bipartite graph of
\citet{bunz2015faster}.

\paragraph{Research gap.} No reviewed study treats an open-route class and
a detour-constrained class as strategic bidders within one truthful
mechanism, and none characterises how far a bid-independent route pool can
be reduced without affecting the VCG outcome.

\paragraph{Design concepts considered.} The review leads to five design
concepts that shape the proposed system:
\begin{enumerate}[label=(\alph*),leftmargin=1cm,itemsep=2pt]
\item \emph{Bid independence}: everything that defines the allocation range
must be computed from public data before reports are collected;
\item \emph{Exactness}: winner determination and counterfactuals must be
solved to proven optimality;
\item \emph{Locality}: pruning decisions for a driver may use only that
driver's own data and public prices, to avoid an NP-hard global reasoning
step;
\item \emph{A priced outside option}: the fixed fleet's public price is the
only information about the rest of the market that a local rule can rely on;
\item \emph{A single private parameter} per driver, which keeps the report
space one-dimensional and the cost of every route affine in the report.
\end{enumerate}

\section{Key References / Candidate Solution Methods}
\label{sec:candidates}

The system has four components: a mechanism (how winners are chosen and
paid), route generation, pool reduction and payment computation. For each
component this section describes the candidate methods found in the key
references or developed during the preliminary rounds of this thesis.

\subsection{Candidate mechanisms}

\begin{itemize}[leftmargin=*,itemsep=3pt]
\item \emph{Posted prices.} The platform announces a price per order,
for example a fraction $\lambda$ of the fixed-fleet price, and drivers
accept routes whose posted payment covers their cost. Simple and
truthful in a weak sense (drivers only accept or decline), but the
allocation can be inefficient and the price must be tuned.
\item \emph{Pay-as-bid (first price).} Winners are paid their reported cost.
Drivers then have an incentive to inflate their reports, and the outcome
depends on how they reason about each other; there is no dominant strategy.
\item \emph{Exact VCG} \citep{zou2022mechanisms,li2026auction}. Exact
optimisation on a fixed range with Clarke-pivot payments; truthful and
individually rational, at the price of one extra exact solve per winner.
\item \emph{Approximate or greedy VCG-like mechanisms}
\citep{li2026auction,wilkens2012single}. Avoid repeated exact solves, but
lose either exact efficiency within the range or deterministic truthfulness.
\end{itemize}

\subsection{Candidate route-generation methods}

\begin{itemize}[leftmargin=*,itemsep=3pt]
\item \emph{Brute-force enumeration}: enumerate every bundle of at most $B$
orders and every precedence-feasible sequence of it; complete but
explosive, and therefore used only as a test oracle.
\item \emph{Level-wise insertion (breadth-first)}: build sequences of
bundles with $k+1$ orders by inserting one order into sequences of bundles
with $k$ orders, checking feasibility with forward slack.
\item \emph{Level-wise insertion with dominance}: as above, but keep only
sequences that are not dominated by another sequence of the same bundle,
using summary attributes (end node, time, cost, slack).
\item \emph{Forward label-setting with exact-key dominance}
\citep{dumas1991pickup,gschwind2018bidirectional}: build routes forward in
time, one stop at a time, and compare only labels that share the current
node, the set of orders on board and the set of delivered orders.
\item \emph{Column generation}: generate only routes with negative reduced
cost at the current dual values. Its pricing step depends on the reported
costs, so the pool would depend on the bids.
\item \emph{Detour-monotone branching} \citep{li2026auction}: discard a
bundle and all its supersets once a lower bound on its detour exceeds the
budget; applicable to occasional drivers only.
\end{itemize}

\subsection{Candidate pool-reduction rules}

\begin{itemize}[leftmargin=*,itemsep=3pt]
\item \emph{No reduction}: run the mechanism on the full pool.
\item \emph{Same-bundle Pareto filter}: for each bundle keep only routes with
Pareto-minimal (distance cost, working time); safe because reports are
nonnegative.
\item \emph{Per-driver hull across bundles}: delete a route if another route
of the same driver is cheaper in both cost components, even if it serves
different orders.
\item \emph{Coarser label keys or bounds} during generation: lookahead
profiles, self-survival bounds, or merging labels with different delivered
sets.
\item \emph{Instance-specific reduction}: keep only routes that are optimal
for some sampled report profile.
\item \emph{Non-local exact reduction}: delete a route if, using the public
data of all drivers, it can be shown never to be needed.
\item \emph{Local FD-completion frontier} (proposed in this thesis): delete
a route if, for every admissible report, some subset route of the same
driver with the remaining orders sent to the fixed fleet is no more
expensive.
\end{itemize}

\subsection{Candidate payment-computation methods}

\begin{itemize}[leftmargin=*,itemsep=3pt]
\item \emph{Naive re-solve}: solve the winner-determination problem once
without each winner, using the same pool and tie-breaking.
\item \emph{Component decomposition}: re-solve only the part of the problem
that contains the removed winner, using a conflict graph built from the
pool \citep[in the spirit of][]{sandholm2005cabob,bunz2015faster}.
\item \emph{Approximate counterfactuals}
\citep{wilkens2012single,li2026auction}: avoid exact re-solves, at the cost of
exact truthfulness.
\end{itemize}
